\documentclass[../../main.tex]{subfiles}

\begin{document}

\section{局部紧拓扑空间}\label{section:ch21-局部紧}

拓扑空间 $X$ 称为是局部紧的, 若 $X$ 中的每个点有紧闭包的邻域. 每个紧空间是局部紧的, 而欧氏空间 $\mathbf{R}^{n}$ 是局部紧的但不是紧的空间的例子. Riesz 定理告诉我们具有由范数诱导的拓扑的无穷维赋范线性空间不是局部紧的. 本节我们证明局部紧空间的性质, 它们是我们接下来对测度与拓扑的研究的基础.

\textbf{Urysohn 引理的变体.} 回忆我们推广邻域这个词的意义, 对拓扑空间 $X$ 的子集 $K$, 称包含 $K$ 的开集为 $K$ 的邻域.

\begin{lemma}\label{lemma:ch21-点邻域紧闭包}
令 $x$ 为局部紧 Hausdorff 空间 $X$ 的点而 $\mathcal{O}$ 是 $x$ 的邻域. 则存在 $x$ 的邻域 $\mathcal{V}$, 它具有包含于 $\mathcal{O}$ 的紧闭包, 即
\begin{align*}
x\in\mathcal{V}\subseteq\overline{\mathcal{V}}\subseteq\mathcal{O}\text{ 且 }\overline{\mathcal{V}}\text{ 是紧的}.
\end{align*}
\end{lemma}

\begin{proof}
令 $\mathcal{U}$ 为 $x$ 的具有紧闭包的邻域. 则拓扑空间 $\overline{\mathcal{U}}$ 是紧 Hausdorff 的, 因此是正规的. 集合 $\mathcal{O}\cap\mathcal{U}$ 是 $x$ 的关于 $\overline{\mathcal{U}}$ 拓扑的一个邻域. 因此由 $\overline{\mathcal{U}}$ 的正规性, 存在 $x$ 的邻域 $\mathcal{V}$, 它有包含于 $\mathcal{O}\cap\mathcal{U}$ 的紧闭包. 这里邻域与闭包都是关于 $\overline{\mathcal{U}}$ 拓扑的. 由于 $\mathcal{O}$ 与 $\mathcal{U}$ 在 $X$ 中是开的, 从子空间拓扑的定义得出 $\mathcal{V}$ 在 $X$ 中是开的且 $\overline{\mathcal{V}}\subseteq\mathcal{O}$, 其中闭包是关于 $X$ 的拓扑.
\end{proof}

\begin{proposition}\label{proposition:ch21-紧集邻域}
令 $K$ 为局部紧 Hausdorff 空间 $X$ 的紧子集而 $\mathcal{O}$ 是 $K$ 的邻域. 则存在 $K$ 的邻域 $\mathcal{V}$, 具有包含于 $\mathcal{O}$ 的紧闭包, 即
\begin{align*}
K\subseteq\mathcal{V}\subseteq\overline{\mathcal{V}}\subseteq\mathcal{O}\text{ 且 }\overline{\mathcal{V}}\text{ 是紧的}.
\end{align*}
\end{proposition}

\begin{proof}
根据前一个引理, 每个点 $x\in K$ 有一个邻域 $\mathcal{N}_{x}$ 具有包含于 $\mathcal{O}$ 的紧闭包. 则 $\{\mathcal{N}_{x}\}_{x\in K}$ 是紧集 $K$ 的开覆盖. 选取 $K$ 的有限子覆盖 $\{\mathcal{N}_{x_{i}}\}_{i=1}^{n}$. 集合 $\mathcal{V}=\bigcup_{i=1}^{n}\mathcal{N}_{x_{i}}$ 是 $K$ 的邻域且
\begin{align*}
\overline{\mathcal{V}}\subseteq\bigcup_{i=1}^{n}\overline{\mathcal{N}}_{x_{i}}\subseteq\mathcal{O}.
\end{align*}
集合 $\bigcup_{i=1}^{n}\overline{\mathcal{N}}_{x_{i}}$ 作为紧集的有限族的并是紧的, 因此 $\overline{\mathcal{V}}$ 也如此, 这是由于它是紧空间的闭子集.
\end{proof}

对拓扑空间 $X$ 上的连续函数 $f$, $f$ 的支撑, 我们记为 $\operatorname{supp}f$, 定义为集合 $\{x\in X\mid f(x)\neq 0\}$ 的闭包, 即
\begin{align*}
\operatorname{supp}f=\overline{\{x\in X\mid f(x)\neq 0\}}.
\end{align*}
我们记由具有紧支撑的连续函数 $f:X\to\mathbf{R}$ 组成的集合为 $C_{c}(X)$. 因此函数属于 $C_{c}(X)$ 当且仅当它是连续的且在紧集外消失.

\begin{proposition}\label{proposition:ch21-urysohn变体}
令 $K$ 为局部紧 Hausdorff 空间 $X$ 的子集而 $\mathcal{O}$ 是 $K$ 的邻域. 则存在属于 $C_{c}(X)$ 的函数 $f$ 使得
\begin{align}\label{align:ch21-urysohn函数}
\text{在 }K\text{ 上 }f=1,\quad\text{在 }X\setminus\mathcal{O}\text{ 上 }f=0,\text{在 }X\text{ 上 }0\leqslant f\leqslant 1.
\end{align}
\end{proposition}

\begin{proof}
根据前一个命题, 存在 $K$ 的邻域 $\mathcal{V}$, 它有包含于 $\mathcal{O}$ 的紧闭包. 由于 $\overline{\mathcal{V}}$ 是紧的与 Hausdorff 的, 它是正规的. 此外, $K$ 与 $\overline{\mathcal{V}}\setminus\mathcal{V}$ 是 $\overline{\mathcal{V}}$ 的不交闭子集. 根据 Urysohn 引理, 存在 $\overline{\mathcal{V}}$ 上的连续实值函数 $f$ 使得
\begin{align*}
\text{在 }K\text{ 上 }f=1,\text{在 }\overline{\mathcal{V}}\setminus\mathcal{V}\text{ 上 }f=0,\text{在 }\overline{\mathcal{V}}\text{ 上 }0\leqslant f\leqslant 1.
\end{align*}
则 $f$ 属于 $C_{c}(X)$ 且有\eqref{align:ch21-urysohn函数}描述的性质.
\end{proof}

回忆拓扑空间的子集称为 $G_{\delta}$ 集, 若它是可数个开集的交.

\begin{corollary}\label{corollary:ch21-gdelta函数}
令 $K$ 为局部紧 Hausdorff 空间 $X$ 的紧 $G_{\delta}$ 子集. 则存在函数 $f\in C_{c}(X)$ 使得
\begin{align*}
K=\{x\in X\mid f(x)=1\}.
\end{align*}
\end{corollary}

\begin{proof}
根据命题\cref{proposition:ch21-紧集邻域}, 存在 $K$ 的具有紧闭包的邻域 $\mathcal{U}$. 由于 $K$ 是 $G_{\delta}$ 集, 存在开集的可数族 $\{\mathcal{O}_{k}\}_{k=1}^{\infty}$, 其交是 $K$.

我们可以假设对所有 $k$, $\mathcal{O}_{k}\subseteq\mathcal{U}$. 根据前一个命题, 对每个 $k$ 存在 $X$ 上的连续实值函数 $f_{k}$ 使得
\begin{align*}
\text{在 }K\text{ 上 }f_{k}=1,\text{在 }X\setminus\overline{\mathcal{O}}_{k}\text{ 上 }f_{k}=0,\quad\text{在 }X\text{ 上 }0\leqslant f_{k}\leqslant 1.
\end{align*}
$X$ 上定义为
\begin{align*}
f=\sum_{k=1}^{\infty}2^{-k}f_{k}
\end{align*}
的函数 $f$ 具有所要的性质.
\end{proof}

\textbf{单位分解.}

\begin{definition}[单位分解]\label{definition:ch21-单位分解}
令 $K$ 为拓扑空间 $X$ 的子集, 它可被开集 $\{\mathcal{O}_{k}\}_{k=1}^{n}$ 覆盖. $X$ 上的连续实值函数族 $\{\varphi_{i}\}_{i=1}^{n}$ 称为是 $K$ 的从属于 $\{\mathcal{O}_{k}\}_{k=1}^{n}$ 的单位分解, 若
\begin{align*}
\text{在 }X\text{ 上对 }1\leqslant i\leqslant n,\quad\operatorname{supp}\varphi_{i}\subseteq\mathcal{O}_{i},\quad 0\leqslant\varphi_{i}\leqslant 1
\end{align*}
且
\begin{align*}
\text{在 }K\text{ 上},\varphi_{1}+\varphi_{2}+\cdots+\varphi_{n}=1.
\end{align*}
\end{definition}

\begin{proposition}\label{proposition:ch21-单位分解存在}
令 $K$ 为局部紧 Hausdorff 空间 $X$ 的紧子集, 而 $\{\mathcal{O}_{k}\}_{k=1}^{n}$ 是 $K$ 的有限开集覆盖. 则存在 $K$ 的从属于该有限覆盖的单位分解 $\{\varphi_{k}\}_{k=1}^{n}$ 且每个 $\varphi_{k}$ 有紧支撑.
\end{proposition}

\begin{proof}
我们首先宣称存在 $K$ 的开覆盖 $\{\mathcal{U}_{k}\}_{k=1}^{n}$ 使得对每个 $k$, $\overline{\mathcal{U}}_{k}$ 是 $\mathcal{O}_{k}$ 的紧子集. 事实上, 应用命题\cref{proposition:ch21-紧集邻域} $n$ 次, 对每个 $x\in K$, 存在 $x$ 的具有紧闭包的邻域 $\mathcal{N}_{x}$, 且使得若 $1\leqslant j\leqslant n$ 而 $x$ 属于 $\mathcal{O}_{j}$ 则 $\overline{\mathcal{N}}_{x}\subseteq\mathcal{O}_{j}$. 开集族 $\{\mathcal{N}_{x}\}_{x\in K}$ 是 $K$ 的覆盖且 $K$ 是紧的. 因此存在有限个点 $\{x_{k}\}_{k=1}^{n}$ 使得 $\{\overline{\mathcal{N}}_{x_{k}}\}_{1\leqslant k\leqslant n}$ 也覆盖 $K$. 对 $1\leqslant k\leqslant n$, 令 $\mathcal{U}_{k}$ 为包含于 $\mathcal{O}_{k}$ 的那些 $\overline{\mathcal{N}}_{x_{j}}$ 的并. 则 $\{\mathcal{U}_{1},\cdots,\mathcal{U}_{n}\}$ 是 $K$ 的开覆盖且对每个 $k$, $\overline{\mathcal{U}}_{k}$ 是 $\mathcal{O}_{k}$ 的紧子集. 由于它是这样的集合的有限并. 我们从命题\cref{proposition:ch21-urysohn变体}推出对每个 $k$($1\leqslant k\leqslant n$), 存在函数 $f_{k}\in C_{c}(X)$, 在 $\overline{\mathcal{U}}_{k}$ 上 $f_{k}=1$ 而在 $X\setminus\mathcal{O}_{k}$ 上 $f=0$. 同一个命题告诉我们存在函数 $h\in C(X)$, 使得在 $K$ 上 $h=1$ 而在 $X\setminus\bigcup_{k=1}^{n}\mathcal{U}_{k}$ 上 $h=0$. 在 $X$ 上定义
\begin{align*}
f=\sum_{k=1}^{n}f_{k}.
\end{align*}
观察到在 $X$ 上 $f+[1-h]>0$ 而在 $K$ 上 $h=0$. 因此若我们在 $X$ 上对 $1\leqslant k\leqslant n$ 定义
\begin{align*}
\varphi_{k}=\frac{f_{k}}{f+[1-h]}.
\end{align*}
则 $\{\varphi_{k}\}_{k=1}^{n}$ 是 $K$ 的从属于 $\{\mathcal{O}_{k}\}_{k=1}^{n}$ 的单位分解且每个 $\varphi_{k}$ 有紧支撑.
\end{proof}

\begin{theorem}[Alexandroff 单点紧化定理]\label{theorem:ch21-alexandroff}
若 $X$ 是局部紧 Hausdorff 空间, 通过把不属于 $X$ 的单点 $\omega$ 加入 $X$, 且定义 $X^{*}$ 中的集合为开的(若它是 $X$ 中的开集或它是 $X$ 中的紧子集的补集), 我们能够建立一个新的空间 $X^{*}$. 则 $X^{*}$ 是紧 Hausdorff 空间, $X$ 到 $X^{*}$ 的恒等映射是 $X$ 与 $X^{*}\sim\{\omega\}$ 的同胚. 空间 $X^{*}$ 称为 $X$ 的 Alexandroff 单点紧化, 而 $\omega$ 常称为 $X^{*}$ 中的无穷远点.
\end{theorem}

以下的 Tietze 延拓定理针对局部紧 Hausdorff 空间的变体的证明很好地说明了 Alexandroff 紧化的用途.

\begin{theorem}\label{theorem:ch21-tietze延拓}
令 $K$ 为局部紧 Hausdorff 空间 $X$ 的紧子集而 $f$ 是 $K$ 上的连续实值函数. 则 $f$ 可延拓为整个 $X$ 上的连续实值函数.
\end{theorem}

\begin{proof}
$X$ 的 Alexandroff 紧化 $X^{*}$ 是紧 Hausdorff 空间. 此外, $K$ 是 $X^{*}$ 的闭子集, 这是由于它的补在 $X^{*}$ 中是开的. 紧 Hausdorff 空间是正规的. 因此我们从 Tietze 延拓定理推出 $f$ 可被延拓为整个 $X^{*}$ 上的连续实值函数. 该延拓在 $X$ 上的限制是 $f$ 到整个 $X$ 的连续延拓.
\end{proof}

\subsection*{习题}

\begin{exercise}\label{exercise:ch21-01}
令 $X$ 为局部紧 Hausdorff 空间, 而 $F$ 是与 $X$ 的每个紧子集有闭的交的集合. 证明 $F$ 是闭的.
\end{exercise}

\begin{exercise}\label{exercise:ch21-02}
关于命题\cref{proposition:ch21-urysohn变体}的证明:
\begin{enumerate}[(i)]
\item 证明 $F$ 与 $\overline{\mathcal{V}}\setminus\mathcal{V}$ 是 $\overline{\mathcal{V}}$ 的闭子集.
\item 证明函数 $f$ 是连续的.
\end{enumerate}
\end{exercise}

\begin{exercise}\label{exercise:ch21-03}
令 $X$ 为局部紧 Hausdorff 空间而 $X^{*}$ 是 $X$ 的 Alexandroff 单点紧化:
\begin{enumerate}[(i)]
\item 证明 $X^{*}$ 的那些或者是 $X$ 的开子集或者是 $X$ 的紧子集的补集的子集是 $X^{*}$ 的拓扑.
\item 证明从 $X$ 到子空间 $X^{*}\sim\{\omega\}$ 的恒等映射是同胚.
\item 证明 $X^{*}$ 是紧与 Hausdorff 的.
\end{enumerate}
\end{exercise}

\begin{exercise}\label{exercise:ch21-04}
证明 $\mathbf{R}^{n}$ 的 Alexandroff 单点紧化同胚于 $n$ 维球面 $S^{n}=\{x\in\mathbf{R}^{n+1}\mid\|x\|=1\}$.
\end{exercise}

\begin{exercise}\label{exercise:ch21-05}
证明具有子空间拓扑的局部紧 Hausdorff 空间的开子集是局部紧的.
\end{exercise}

\begin{exercise}\label{exercise:ch21-06}
证明具有子空间拓扑的局部紧空间的闭子集是局部紧的.
\end{exercise}

\begin{exercise}\label{exercise:ch21-07}
证明局部紧 Hausdorff 空间 $X$ 是紧的当且仅当由无穷远点构成的集合是 $X$ 的 Alexandroff 单点紧化 $X^{*}$ 的开子集.
\end{exercise}

\begin{exercise}\label{exercise:ch21-08}
令 $X$ 为局部紧 Hausdorff 空间. 证明 Alexandroff 单点紧化 $X^{*}$ 是可分的当且仅当 $X$ 是可分的.
\end{exercise}

\begin{exercise}\label{exercise:ch21-09}
考虑由实数集组成的拓扑空间 $X$. 它的拓扑以可数集的补集作为基. 证明 $X$ 不是局部紧的.
\end{exercise}

\begin{exercise}\label{exercise:ch21-10}
将 Urysohn 引理用于 $X$ 的 Alexandroff 单点紧化给出命题\cref{proposition:ch21-urysohn变体}的证明.
\end{exercise}

\begin{exercise}\label{exercise:ch21-11}
令 $f$ 为将局部紧 Hausdorff 空间 $X$ 映上拓扑空间 $Y$ 的连续映射. $Y$ 必须是局部紧的吗?
\end{exercise}

\begin{exercise}\label{exercise:ch21-12}
令 $X$ 为拓扑空间, $f$ 是 $X$ 的有紧支撑的连续函数. 定义 $K=\{x\in X\mid f(x)=1\}$. 证明 $K$ 是紧 $G_{\delta}$ 集.
\end{exercise}

\begin{exercise}\label{exercise:ch21-13}
令 $\mathcal{O}$ 为紧 Hausdorff 空间 $X$ 的开子集. 证明 $X$ 到 $\mathcal{O}$ 的 Alexandroff 单点紧化的映射是连续的, 它在 $\mathcal{O}$ 上是恒等映射而将 $X\setminus\mathcal{O}$ 的每一点映到 $\omega$.
\end{exercise}

\begin{exercise}\label{exercise:ch21-14}
令 $X$ 和 $Y$ 为局部紧 Hausdorff 空间, 而 $f$ 是 $X$ 到 $Y$ 的连续映射. 令 $X^{*}$ 和 $Y^{*}$ 为 $X$ 和 $Y$ 的 Alexandroff 单点紧化, 而 $f^{*}$ 是从 $X^{*}$ 映上 $Y^{*}$ 的映射, 其限制在 $X$ 上是 $f$ 且将 $X^{*}$ 的无穷远点映到 $Y^{*}$ 的无穷远点. 证明 $f^{*}$ 是连续的当且仅当只要 $K\subseteq Y$ 是紧的, $f^{-1}(K)$ 就是紧的. 具有该性质的映射 $f$ 称为是适当的.
\end{exercise}

\begin{exercise}\label{exercise:ch21-15}
令 $X$ 为局部紧 Hausdorff 空间. 证明 $X$ 的子集 $F$ 是闭的当且仅当对 $X$ 的每个紧子集 $K$, $F\cap K$ 是闭的. 此外, 证明: 若空间是第一可数而非局部紧, 则同样的等价关系也成立.
\end{exercise}

\begin{exercise}\label{exercise:ch21-16}
令 $\mathcal{F}$ 为局部紧 Hausdorff 空间 $X$ 上的实值连续函数族. 它有以下性质:
\begin{enumerate}[(i)]
\item 若 $f\in\mathcal{F}$ 且 $g\in\mathcal{F}$, 则 $f+g\in\mathcal{F}$.
\item 若 $f\in\mathcal{F}$ 且 $g\in\mathcal{F}$, 则 $f/g\in\mathcal{F}$, 假定 $\operatorname{supp}f\subseteq\{x\in X\mid g(x)\neq 0\}$.
\item 给定点 $x_{0}\in X$ 的邻域 $\mathcal{O}$, 存在 $f\in\mathcal{F}$ 满足 $f(x_{0})=1$, $0\leqslant f\leqslant 1$ 且 $\operatorname{supp}f\subseteq\mathcal{O}$.
\end{enumerate}
证明命题\cref{proposition:ch21-单位分解存在}仍然成立, 若我们要求单位分解中的函数属于 $\mathcal{F}$.
\end{exercise}

\begin{exercise}\label{exercise:ch21-17}
令 $K$ 为局部紧 Hausdorff 空间 $X$ 的紧 $G_{\delta}$ 子集. 证明存在 $X$ 上的递减的连续非负实值函数序列, 它在 $X$ 上逐点地收敛到 $K$ 的特征函数.
\end{exercise}

\begin{exercise}\label{exercise:ch21-18}
Baire 范畴定理断言完备度量空间开稠密子集的可数族的交也是稠密的. 它的证明核心在于 Cantor 交定理. 证明在给出以下断言的证明时, Fr\'echet 交定理是 Cantor 交定理的足够强的替代. 可先对 $X$ 是紧的情形证明: 令 $X$ 为局部紧 Hausdorff 空间.
\begin{enumerate}[(i)]
\item 若 $\{F_{n}\}_{n=1}^{\infty}$ 是 $X$ 的闭子集的可数族, 使得每个 $F_{n}$ 有空的内部, 则并 $\bigcup_{n=1}^{\infty}F_{n}$ 也有空的内部.
\item 若 $\{\mathcal{O}_{n}\}_{n=1}^{\infty}$ 是 $X$ 的开稠密集的可数族, 则交 $\bigcap_{n=1}^{\infty}\mathcal{O}_{n}$ 也是稠密的.
\end{enumerate}
\end{exercise}

\begin{exercise}\label{exercise:ch21-19}
用前一个习题证明: 令 $X$ 为局部紧 Hausdorff 空间. 若 $\mathcal{O}$ 是 $X$ 的开子集, 它包含于 $X$ 的闭子集的可数并 $\bigcup_{n=1}^{\infty}F_{n}$, 则它们的内部的并 $\bigcup_{n=1}^{\infty}\operatorname{int}F_{n}$ 是 $\mathcal{O}$ 的开稠密集.
\end{exercise}

\begin{exercise}\label{exercise:ch21-20}
对映射 $f:X\to Y$ 与 $Y$ 的子集族 $\mathcal{C}$ 定义 $X$ 的子集族 $f^{*}\mathcal{C}$ 为
\begin{align*}
f^{*}\mathcal{C}=\{E\mid E=f^{-1}[C],\ C\in\mathcal{C}\}.
\end{align*}
证明: 若 $\mathcal{A}$ 是由 $\mathcal{C}$ 生成的 $\sigma$ 代数, 则 $f^{*}\mathcal{A}$ 是由 $f^{*}\mathcal{C}$ 生成的 $\sigma$ 代数.
\end{exercise}

\begin{exercise}\label{exercise:ch21-21}
对映射 $f:X\to Y$ 与 $X$ 的子集族 $\mathcal{C}$, 令 $\mathcal{A}$ 为 $\mathcal{C}$ 生成的 $\sigma$ 代数. 若对每个 $C\in\mathcal{C}$, $f^{-1}[f[C]]=C$, 证明对每个 $A\in\mathcal{A}$, $f^{-1}[f[A]]=A$.
\end{exercise}

\end{document}